\section{Lean interface and trust boundary}
\label{app:lean-interface}

This appendix is intended to make the Lean formalization auditable through a
human mathematical interface.  It records exact Lean declarations rather than
merely matching names.  The present subsection covers only the public inputs
of the manifold-level invariant.  Later subsections must separately flatten
the category laws, topology axioms, derived proof obligations, public outputs,
and final axiom closures.

The source snapshot inspected here is commit
\texttt{aa5ce6c5713dc95e3f80e748d3ea245e0bfdbb8f}.  A source-only red-team scan
of
\nolinkurl{lean/TwistedYetter/Kirby/Final/Semantics/Independent}
reported $19$ explicit axioms and $33$ possible proof fields.  These counts are
conservative syntactic findings, not the axiom closure of the final theorem and
not a judgment that every field is improper.  No kernel audit is claimed in
this subsection.  Its current terminal status is
\texttt{INCONCLUSIVE\_HUMAN\_REVIEW\_REQUIRED}.

\subsection{Public final-theorem inputs}
\label{app:lean-public-inputs}

\begin{enumerate}[label=\textbf{I\arabic*.},leftmargin=*]

\item \textbf{Coefficient field.}

\emph{Human statement.}
The Lean construction is parametrized by an algebraically closed field $K$ of
characteristic zero.  The mathematical results in this paper specialize to
$K=\mathbb C$.

\emph{Lean declaration.}
Every public declaration carries the ambient parameters
\begin{verbatim}
{K : Type uK} [Field K] [IsAlgClosed K] [CharZero K]
\end{verbatim}

\emph{Classification and judgment.}
These are typeclass assumptions, not project-specific axioms.  The Lean
statement is more general than the paper statement.  Thus this is a faithful
specialization at $K=\mathbb C$, but not a literal statement match.

\item \textbf{Twisted premodular category.}

\emph{Human statement.}
The algebraic input is a premodular fusion category $\cC$ over $K$, together
with a monoidal natural automorphism
$\eta\colon\operatorname{Id}_{\cC}\Rightarrow
\operatorname{Id}_{\cC}$.

\emph{Lean declaration.}
\begin{verbatim}
structure TwistedPremodularCategory
    (K : Type uK) [Field K] [IsAlgClosed K] [CharZero K] where
  category : CategoryTheory.BundledPremodularCategory.{v, uK, uC, uI} K
  character : TensorIdentityCharacter (C := category.C)
\end{verbatim}
The declaration occurs in
\nolinkurl{lean/TwistedYetter/CategoryTheory/Premodular/Bundled.lean},
lines 446--449.

\emph{Classification and judgment.}
This is a two-field data structure.  Its immediate human meaning agrees with
the paper.  This does not yet certify the claim: the inherited laws inside
\texttt{BundledPremodularCategory} and \texttt{TensorIdentityCharacter} must be
flattened and reviewed in the category-law subsection.

\item \textbf{Four-manifold vocabulary.}

\emph{Human statement.}
The intended input $X$ is a closed, connected, oriented smooth
four-manifold, and it has an underlying topological space.

\emph{Lean declarations.}
\begin{verbatim}
axiom ClosedConnectedOrientedSmoothFourManifold : Type

axiom underlyingTopologicalSpace :
  ClosedConnectedOrientedSmoothFourManifold → TopCat
\end{verbatim}
They occur in \texttt{EntranceAxioms.lean}, lines 332--336, under
\nolinkurl{lean/TwistedYetter/Kirby/Final/Semantics/Independent}.

\emph{Classification and judgment.}
These are explicit foundational vocabulary axioms; Lean does not construct
smooth four-manifolds from more primitive objects.  Connectedness is now part
of the primitive type's stated meaning and name, rather than being inferred
indirectly from a later presentation-existence axiom.  This exactly matches
the manifold scope of the paper.  Status: \texttt{MATCH}.

\item \textbf{Oriented diffeomorphism.}

\emph{Human statement.}
For four-manifolds $X$ and $Y$, an input $f$ is an orientation-preserving
diffeomorphism $X\to Y$.  Forgetting smooth structure gives a homeomorphism of
their underlying topological spaces.

\emph{Lean declarations.}
\begin{verbatim}
axiom OrientedDiffeomorphism :
  ClosedConnectedOrientedSmoothFourManifold →
  ClosedConnectedOrientedSmoothFourManifold →
  Type

axiom underlyingHomeomorphism
    {X Y : ClosedConnectedOrientedSmoothFourManifold} :
    OrientedDiffeomorphism X Y →
    (underlyingTopologicalSpace X ≅ underlyingTopologicalSpace Y)
\end{verbatim}
The source is \nolinkurl{EntranceAxioms.lean}, lines 338--347.

\emph{Classification and judgment.}
Both are explicit foundational axioms.  Their intended meaning matches the
paper at the level currently stated, but groupoid operations and orientation
compatibility are not visible in these declarations.  Whether downstream
proofs require more than the displayed forgetting map remains part of the
topology-boundary audit.

\item \textbf{Integral second homology and its transport.}

\emph{Human statement.}
For a four-manifold $X$,
\[
  \operatorname{ManifoldH}_2(X)=H_2(X;\mathbb Z).
\]
An oriented diffeomorphism $f\colon X\to Y$ induces the usual
additive-group isomorphism
\[
  f_*\colon H_2(X;\mathbb Z)\longrightarrow H_2(Y;\mathbb Z).
\]

\emph{Lean declarations.}
\begin{verbatim}
noncomputable def ManifoldH₂
    (X : ClosedConnectedOrientedSmoothFourManifold) : AddCommGrpCat :=
  integralSingularH₂Functor.obj (underlyingTopologicalSpace X)

noncomputable def inducedH₂
    {X Y : ClosedConnectedOrientedSmoothFourManifold}
    (f : OrientedDiffeomorphism X Y) :
    ManifoldH₂ X ≃+ ManifoldH₂ Y :=
  CategoryTheory.Iso.addCommGroupIsoToAddEquiv
    (integralSingularH₂Functor.mapIso (underlyingHomeomorphism f))
\end{verbatim}
The source is \nolinkurl{EntranceAxioms.lean}, lines 350--366.

\emph{Classification and judgment.}
These are derived noncomputable definitions, not new axioms.  They use
Mathlib's integral singular-homology functor and the explicitly supplied
underlying homeomorphism.  This matches the paper's $H_2(X;\mathbb Z)$ input.

\item \textbf{The marked public input.}

\emph{Human statement.}
The manifold-side argument of the invariant is exactly a pair
$(X,\alpha)$ with
$\alpha\in H_2(X;\mathbb Z)$; no Kirby diagram, evaluator, compiler, or proof
certificate is supplied by the caller.

\emph{Lean declaration as used by the public definition.}
\begin{verbatim}
(Xα : Σ X : ClosedConnectedOrientedSmoothFourManifold, ManifoldH₂ X)
\end{verbatim}
This is the dependent sum used in
\nolinkurl{lean/TwistedYetter/Kirby/Final/TwistedYetter.lean}, lines 173--177.

\emph{Classification and judgment.}
This is ordinary dependent-pair data, not an axiom or proof field.  It matches
the paper's marked input.

\end{enumerate}

\subsection{Transitive algebraic input}
\label{app:lean-category-input}

This subsection expands the category-theoretic content hidden behind I2.  The
purpose is not to reproduce Mathlib's definitions of category, abelian
category, monoidal category, braiding, or rigidity.  It is to expose every
project-defined layer between those standard interfaces and the public pair
$(\cC,\eta)$, so that no evaluator-correctness statement can hide inside the
word ``premodular.''

\begin{enumerate}[label=\textbf{A\arabic*.},leftmargin=*]

\item \textbf{The first-class bundle.}

\emph{Human statement.}
The object $\cC$ packages one category together with compatible instances
making it abelian, $K$-linear, monoidal, rigid, braided, monoidally additive,
monoidally $K$-linear, and premodular.

\emph{Lean declaration.}
The structure \texttt{BundledPremodularCategory} has precisely the fields
\begin{verbatim}
C, category, abelian, linear, monoidal,
rigid, braided, monoidalPreadditive,
monoidalLinear, premodular
\end{verbatim}
Its complete declaration is in
\nolinkurl{lean/TwistedYetter/CategoryTheory/Premodular/Bundled.lean}, lines
19--37.

\emph{Classification and judgment.}
The first nine fields install standard Mathlib structures and their coherence
laws.  The last field is expanded in A2.  None mentions Kirby diagrams,
graphical evaluation, homology markings, or invariance.  Status:
\texttt{MATCH}, subject to A2--A6.

\item \textbf{Premodularity.}

\emph{Human statement.}
A premodular category over $K$ is a braided ribbon fusion category over an
algebraically closed field of characteristic zero.

\emph{Lean declaration.}
After listing the ambient Mathlib instances, the complete project-specific
class body is
\begin{verbatim}
class PremodularCategory ...
    extends FusionCategory K C, RibbonCategory C
\end{verbatim}
in \nolinkurl{lean/TwistedYetter/CategoryTheory/Premodular/Basic.lean}, lines
25--31.

\emph{Classification and judgment.}
This is a domain-defining class, not a theorem assumption about the proposed
invariant.  It agrees with the paper's use of a premodular fusion category and
with the convention of \cite[Definition~8.13.1]{EGNO2015}.  Its two parents are
expanded next.  Status: \texttt{MATCH}.

\item \textbf{Fusion structure.}

\emph{Human statement.}
A fusion category is a finite semisimple $K$-linear rigid monoidal category
whose tensor unit is simple.

\emph{Lean declaration.}
The class \texttt{FusionCategory} extends
\texttt{FiniteSemisimpleCategory K C} and has the single additional field
\begin{verbatim}
tensorUnit_simple : Simple (tensor unit)
\end{verbatim}
in \nolinkurl{lean/TwistedYetter/CategoryTheory/Fusion/Basic.lean}, lines
25--32.  Here ``tensor unit'' is the source expression \texttt{(one\_ C)}
written without its Unicode glyph.

\emph{Classification and judgment.}
Simplicity of the tensor unit is part of the definition of a fusion category,
not a project-specific proof shortcut.  The class assumes the standard
Mathlib linear, rigid, and monoidal interfaces rather than redefining their
coherence equations.  Status: \texttt{MATCH}.

\item \textbf{Finite semisimplicity and its witnesses.}

\emph{Human statement.}
Every morphism space is finite-dimensional over $K$; there are finitely many
isomorphism classes of simple objects; and every object is a finite direct sum
of simple objects.

\emph{Lean declaration.}
The class \texttt{FiniteSemisimpleCategory} has exactly three fields:
\begin{verbatim}
hom_finite
finite_simple_skeleton
finite_simple_decomposition
\end{verbatim}
in
\nolinkurl{lean/TwistedYetter/CategoryTheory/FiniteSemisimple/Basic.lean},
lines 20--26.  The skeleton witness consists of a finite index type, a family
of simple objects, pairwise nonisomorphism, and completeness among all simple
objects.  For every object $X$, the decomposition witness supplies finitely
many simple summands with inclusions and projections satisfying
\[
 p_i\iota_i=\operatorname{id},
 \qquad p_j\iota_i=0\quad(i\ne j),
 \qquad \sum_i\iota_i p_i=\operatorname{id}_X,
\]
with the composition order translated from Lean's diagrammatic convention.
These witness structures are declared in
\nolinkurl{lean/TwistedYetter/CategoryTheory/FiniteSemisimple/Chosen.lean},
lines 17--52.

\emph{Classification and judgment.}
The displayed equations are the ordinary finite-biproduct coordinate laws.
They are proof fields because they certify that the supplied inclusions and
projections really form a simple decomposition; they do not assert any
Twisted Yetter calculation.  The public caller does not choose these
coordinates: Lean obtains them noncomputably from the existential class
fields.  Status: \texttt{MATCH}.

\item \textbf{Ribbon structure.}

\emph{Human statement.}
A balancing is a natural automorphism
$\theta\colon\operatorname{Id}_{\cC}\Rightarrow
\operatorname{Id}_{\cC}$ satisfying
\[
 \theta_{X\otimes Y}
 =(\theta_X\otimes\theta_Y)\circ c_{Y,X}\circ c_{X,Y}.
\]
It is ribbon when it is compatible with right duality:
\[
 (\theta_X)^*=\theta_{X^*}.
\]

\emph{Lean declaration.}
The structure \texttt{Balancing} contains the natural isomorphism
\texttt{twist} and the displayed tensor equation.  The structure
\texttt{RibbonStructure} extends it by the displayed duality equation, and
\texttt{RibbonCategory} contains one such ribbon structure.  The declarations
are in \nolinkurl{lean/TwistedYetter/CategoryTheory/Ribbon/Basic.lean}, lines
25--49.

\emph{Classification and judgment.}
These are precisely domain-defining ribbon laws.  Pivotal and spherical
structures are not supplied as independent fields: they are constructed from
the ribbon structure, with the spherical identity proved in
\texttt{Ribbon/CanonicalSpherical.lean}.  This prevents callers from supplying
an incompatible pivotal structure.  Status: \texttt{MATCH}; the cited
convention is \cite[Definition~8.10.1 and Proposition~8.10.12]{EGNO2015}.

\item \textbf{The character.}

\emph{Human statement.}
The twist parameter is a monoidal natural automorphism
\[
 \eta\colon\operatorname{Id}_{\cC}\Longrightarrow
 \operatorname{Id}_{\cC}.
\]
Thus every object $X$ has an invertible endomorphism $\eta_X$, natural in
$X$, with
\[
 \eta_{X\otimes Y}=\eta_X\otimes\eta_Y,
 \qquad \eta_{\mathbf 1}=\operatorname{id}_{\mathbf 1}.
\]

\emph{Lean declaration.}
\begin{verbatim}
structure TensorIdentityCharacter where
  iso : identity functor is naturally isomorphic to itself
  monoidal : the forward natural transformation is monoidal
\end{verbatim}
This is a notation-normalized transcription of the complete two-field
structure in
\nolinkurl{lean/TwistedYetter/CategoryTheory/Monoidal/IdentityCharacter.lean},
lines 19--22.

\emph{Classification and judgment.}
Naturality and invertibility are carried by the natural isomorphism;
monoidality supplies the tensor and unit equations.  The action on duals and
the scalar action on simple objects are proved later as theorems, rather than
being extra character fields.  No universal-grading-group presentation is
assumed.  Status: \texttt{MATCH}.

\item \textbf{Derived choices and dimensions.}

\emph{Human statement.}
After receiving $\cC$, Lean may noncomputably choose representatives of simple
objects and decompositions.  Quantum dimensions are canonical ribbon traces,
and the global dimension is
\[
 \operatorname{Dim}(\cC)
 =\sum_{i\in\operatorname{Irr}(\cC)} d_i^2.
\]

\emph{Lean declaration.}
The definitions \texttt{simpleSkeleton}, \texttt{decomposition},
\texttt{chosenSimpleFamily}, \texttt{quantumDimension}, and
\texttt{globalDimension} occur in
\nolinkurl{lean/TwistedYetter/CategoryTheory/Premodular/Bundled.lean}.  They are
all definitions or noncomputable choices derived from A3--A5; none is a field
of \texttt{TwistedPremodularCategory}.

\emph{Classification and judgment.}
These choices affect coordinates used to write finite sums, but not the
public algebraic input.  The only additional literature axiom used by the
final route is the nonvanishing of global dimension; it will be listed in the
separate explicit-axiom subsection.  Status:
\texttt{MATCH\_WITH\_EXPLICIT\_NONVANISHING\_BOUNDARY}.

\end{enumerate}

The transitive project-defined algebraic input therefore contains no
evaluator equality, Kirby-move invariance, representative independence, or
four-manifold realization claim.  The proof fields exposed above define the
standard algebraic domain itself.  The next audit unit is the combinatorial
special Kirby diagram supplied to the evaluator.
